\documentclass[aspectratio=169,11pt]{beamer}
\usepackage[T1]{fontenc}
\usepackage[utf8]{inputenc}
\usepackage[french]{babel}
\usepackage{lmodern}
\usepackage{booktabs}
\usepackage{tabularx}
\usepackage{tikz}
\usetikzlibrary{positioning,arrows.meta,calc}

\usetheme[progressbar=frametitle,numbering=fraction]{metropolis}

\definecolor{nuit}{HTML}{14213D}
\definecolor{papier}{HTML}{F7F6F2}
\definecolor{bleu}{HTML}{1F5FAD}
\definecolor{bleuclair}{HTML}{DCE7F6}
\definecolor{orange}{HTML}{C2581F}
\definecolor{orangeclair}{HTML}{F6E6DA}
\definecolor{vert}{HTML}{2E7D32}
\definecolor{gris}{HTML}{4A5568}
\definecolor{trait}{HTML}{DDD8CC}

\setbeamercolor{normal text}{fg=nuit,bg=papier}
\setbeamercolor{frametitle}{fg=papier,bg=nuit}
\setbeamercolor{progress bar}{fg=bleu,bg=bleuclair}
\setbeamercolor{alerted text}{fg=orange}
\setbeamercolor{title separator}{fg=bleu}
\setbeamercolor{background canvas}{bg=papier}

\tikzset{
  carte/.style={draw=trait,fill=white,rounded corners=4pt,align=center,inner sep=6pt,font=\small},
  fleche/.style={-{Stealth[length=2.5mm]},thick,bleu},
}

\title{Dessiner un graphe directement en espace écran}
\subtitle{Notre proposition d'orientation et nos questions}
\author{Melissa Colin \and Quentin Lopez \and Yani Haouchine}
\institute{ENSEIRB-MATMECA 3A, projet semestriel \quad·\quad Encadrant : Romain Giot (LaBRI)}
\date{Octobre 2026}

\begin{document}

\maketitle

%---------------------------------------------------------------------
\begin{frame}{Notre idée : un réseau qui place les nœuds directement en pixels}
\centering
\begin{tikzpicture}[node distance=0.32cm]
  \node[carte] (g) {\textbf{Graphe}\\\scriptsize nœuds, arêtes, tailles};
  \node[carte,right=of g,fill=bleuclair,draw=bleu] (gnn) {\textbf{Modèle}\\\scriptsize ex. GNN};
  \node[carte,right=of gnn] (pos) {\textbf{Positions}\\\scriptsize $(N,2)$ en pixels};
  \node[carte,right=of pos,fill=orangeclair,draw=orange] (cout) {\textbf{Loss}\\\scriptsize ex. stress, chevauchement};
  \draw[fleche] (g) -- (gnn);
  \draw[fleche] (gnn) -- (pos);
  \draw[fleche] (pos) -- (cout);
  \draw[-{Stealth[length=2.5mm]},thick,orange,dashed] (cout.south) -- ++(0,-0.55) -| node[pos=0.25,below,font=\scriptsize]{le gradient remonte jusqu'aux poids du modèle} (gnn.south);
\end{tikzpicture}
\end{frame}

%---------------------------------------------------------------------
\begin{frame}{Soft-argmax : une sortie dans l'espace des pixels}
\centering
\begin{tikzpicture}
  % ---------------- gauche : origine, estimation de pose humaine
  \node[font=\small\bfseries] at (1.6,3.35) {Origine : estimation de pose humaine};
  \begin{scope}[shift={(1.2,0.25)},line width=1.6pt,nuit,line cap=round]
    \draw (0,2.25) circle (0.25);
    \draw (0,2.0) -- (0,1.0);
    \draw (0,1.75) -- (-0.45,1.25) -- (-0.65,0.75);
    \draw (0,1.75) -- (0.5,1.35) -- (0.85,1.7);
    \draw (0,1.0) -- (-0.35,0.15);
    \draw (0,1.0) -- (0.35,0.15);
    \fill[orange!18] (0.85,1.7) circle (0.42);
    \fill[orange!40] (0.85,1.7) circle (0.26);
    \fill[orange!85] (0.85,1.7) circle (0.1);
  \end{scope}
  \draw[-{Stealth[length=2mm]},thick,orange] (2.35,1.95) -- (3.0,1.95) node[right,font=\scriptsize,nuit,align=left] {$(x, y)$\\du poignet};
  \node[font=\scriptsize,gris,align=center] at (1.6,-0.15) {carte de chaleur d'une articulation\\$\rightarrow$ sa position dans l'image};

  % ---------------- séparation
  \draw[trait,thick] (4.25,-0.4) -- (4.25,3.5);

  % ---------------- droite : chez nous, un nœud
  \node[font=\small\bfseries] at (8.1,3.35) {Chez nous : un nœud du graphe};
  \begin{scope}[shift={(5.1,0.15)},scale=0.36]
    \foreach \i in {0,...,7} {
      \foreach \j in {0,...,7} {
        \pgfmathsetmacro{\v}{90*exp(-((\i-4.4)^2+(\j-3.3)^2)/2.6)}
        \fill[bleu!\v!white] (\i,\j) rectangle ++(1,1);
      }
    }
    \draw[step=1,trait] (0,0) grid (8,8);
    % position moyenne (continue, entre deux cases)
    \fill[orange] (4.95,3.85) circle (0.28);
    \draw[orange,thick] (4.95,3.85) circle (0.5);
  \end{scope}
  \draw[-{Stealth[length=2mm]},thick,orange] (7.2,1.55) -- (8.25,1.55) node[right,font=\scriptsize,nuit,align=left] {position $(x, y)$\\continue,\\entre deux pixels};
  \node[font=\scriptsize,gris,align=center] at (6.65,-0.15) {probabilité sur chaque pixel\\(softmax)};

  % ---------------- formule
  \node[font=\small] at (4.55,-1.15) {$\displaystyle (x, y) \;=\; \sum_{\text{pixels } p} \mathrm{softmax}(s)_p \cdot \mathrm{coord}(p)$};
  \node[font=\scriptsize,gris] at (4.55,-1.75) {moyenne pondérée : différentiable, contrairement à argmax};
  \node[font=\tiny,gris] at (4.55,-2.15) {Integral Pose Regression (Sun et al., 2018) \quad·\quad DSNT (Nibali et al., 2018)};
\end{tikzpicture}
\end{frame}

%---------------------------------------------------------------------
\begin{frame}{Entraînement : sans dessins de référence, en trois phases}
\centering
\begin{tikzpicture}[node distance=0.45cm]
  \node[carte,text width=3.9cm,draw=bleu,line width=0.8pt] (p1) {\textcolor{bleu}{\scriptsize PHASE 1}\\\textbf{Préentraînement}\\[2pt]\scriptsize Le stress seul, calculé à partir des distances dans le graphe : une structure globale correcte.};
  \node[carte,text width=3.9cm,draw=vert,line width=0.8pt,right=of p1] (p2) {\textcolor{vert}{\scriptsize PHASE 2 · CŒUR DU SUJET}\\\textbf{Coûts visuels}\\[2pt]\scriptsize Chevauchement, occlusion arête-nœud et densité.};
  \node[carte,text width=3.9cm,draw=orange,line width=0.8pt,right=of p2] (p3) {\textcolor{orange}{\scriptsize PHASE 3 · BONUS}\\\textbf{Préférences}\\[2pt]\scriptsize Un critique appris sur nos comparaisons, puis un affinage par RL, DPO ou rétropropagation directe.};
  \draw[fleche] (p1) -- (p2);
  \draw[fleche] (p2) -- (p3);
\end{tikzpicture}
\end{frame}

%---------------------------------------------------------------------
\begin{frame}{Une piste à tester : rendre le modèle attentif aux arêtes}
\begin{columns}[T]
\begin{column}{0.48\textwidth}
\begin{block}{A. Tête auxiliaire}
\small Une tête en plus prédit les défauts de chaque arête (croisements, nœud caché, longueur). Entraînée par une perte auxiliaire, elle force le modèle à encoder l'information sur les arêtes dans sa représentation.
\end{block}
\end{column}
\begin{column}{0.48\textwidth}
\begin{block}{B. Raffinement itératif}
\small Une première passe propose des positions. On calcule exactement les défauts de chaque arête, puis une seconde passe les reçoit dans son attention et corrige.
\end{block}
\end{column}
\end{columns}
\end{frame}

%---------------------------------------------------------------------
\begin{frame}{Nos questions}
\begin{enumerate}
\setlength{\itemsep}{8pt}
\item Rôle des SDF : pour quelles losses un rendu différentiable est-il nécessaire, face à une loss géométrique ?
\item Entrées de départ : structure et taille des nœuds seulement, ou aussi épaisseurs, couleurs, formes ?
\item Ordres de grandeur visés : taille des graphes, résolution d'image, temps de calcul
\item Méthodes de référence prioritaires, et implémentations à réutiliser
\item Attendus des jalons du 06/11 et du 18/12 : état de l'art écrit ou oral ?
\item Accès aux exemples de code mentionnés dans le sujet
\end{enumerate}
\end{frame}

\end{document}
